
\chapter{Summary}

\section{Contributions and conclusions}

The thesis developed three LLM-based methods for constructing semantic triples,
a QLoRA procedure for distilling the conversion process into domain-specialized
models, and an OpenEvolve-based method for constructing executable
surface-realization programs. Together, these components provide explicit
meaning representations and reusable realization code rather than a single
opaque generation step.

Text-First Extraction with Normalization is the strongest reported pipeline.
Its post-hoc predicate normalization provides broad semantic coverage, whereas
Evolving Catalog and Blueprint-Iterative Refinement provide stronger up-front
vocabulary control and more directly auditable artifacts. Fine-tuning improves
triple extraction substantially over the base model, but the results remain
dependent on the domain and the selected evaluation metric. The fine-tuned model
is an efficient distillation of pipeline-produced targets rather than
independent semantic supervision.

For surface realization, evolutionary program construction produces an
inspectable Python artifact that can be executed without an inference-time LLM.
The experiments show that the evolutionary fitness function affects the quality
that is optimized: reference-based and LLM-judge objectives favor different
properties. The evolved programs do not outperform fine-tuned BART or direct
prompting on every reported metric. Their principal advantage is transparency,
reusability, and the ability to inspect or modify the realization logic.

\section{Limitations and future work}

Automatic and human evaluation should be used together. The observed saturation
of some LLM-judge scores and the limited single-rater human evaluation show that
neither perspective alone establishes factual correctness. The complete
whole-system evaluation has not been done; consequently, this thesis
does not claim measured end-to-end performance. Its future completion should be
supplemented with independently annotated examples, multiple human raters,
calibrated judges, and comparisons of latency, cost, and maintainability.

In conclusion, this work establishes a practical framework for interpretable
data-to-text generation. The current results favor flexible text-first
extraction with post-hoc normalization for meaning representation and position
evolved programs as transparent alternatives to direct prompting. The remaining
end-to-end evaluation could determine how these component-level advantages
transfer to complete generation from raw structured data.
